% !TEX program = xelatex
% Self-contained source. Compile with XeLaTeX twice.
% Assignment 3 - Bond futures and cheapest-to-deliver analysis.
\documentclass[11pt]{article}
\usepackage[a4paper,top=1in,bottom=1in,left=1.1in,right=1.1in]{geometry}
\usepackage{setspace}
\setstretch{1.18}
\setlength{\parindent}{1.5em}
\setlength{\parskip}{2pt}
\setlength{\emergencystretch}{1.5em}
\usepackage{fontspec}
\setmainfont{texgyreheros-regular.otf}[
  BoldFont=texgyreheros-bold.otf,
  ItalicFont=texgyreheros-italic.otf,
  BoldItalicFont=texgyreheros-bolditalic.otf]
\setsansfont{texgyreheros-regular.otf}[BoldFont=texgyreheros-bold.otf]
\usepackage{amsmath,amssymb,booktabs,array,tabularx}
\usepackage{enumitem,xcolor,colortbl,titlesec,fancyhdr}
\usepackage{float,caption}
\usepackage[hidelinks]{hyperref}

\definecolor{NYUviolet}{HTML}{57068C}
\definecolor{NYUviolet2}{HTML}{8900E1}
\definecolor{inkgrey}{HTML}{3A3A3A}
\definecolor{rulegrey}{HTML}{B9B9B9}
\definecolor{tablehead}{HTML}{ECE3F4}
\definecolor{lightviolet}{HTML}{F7F1FB}
\definecolor{softred}{HTML}{FCEBEC}
\definecolor{softgreen}{HTML}{EAF5EE}

\hypersetup{
  colorlinks=true,
  linkcolor=NYUviolet,
  urlcolor=NYUviolet2,
  citecolor=NYUviolet,
  pdftitle={Assignment 3: Bond Futures and Cheapest-to-Deliver Analysis},
  pdfauthor={Ziyi Yang}}

\titleformat{\section}{\sffamily\Large\bfseries\color{NYUviolet}}
  {\thesection}{0.7em}{}[\vspace{2pt}{\color{rulegrey}\titlerule[0.8pt]}]
\titleformat{\subsection}{\sffamily\large\bfseries\color{inkgrey}}
  {\thesubsection}{0.6em}{}
\titlespacing*{\section}{0pt}{16pt}{8pt}
\titlespacing*{\subsection}{0pt}{10pt}{5pt}

\pagestyle{fancy}
\fancyhf{}
\setlength{\headheight}{15pt}
\renewcommand{\headrulewidth}{0.6pt}
\fancyhead[L]{\footnotesize\sffamily\color{inkgrey}Bond Futures and CTD}
\fancyhead[R]{\footnotesize\sffamily\color{inkgrey}FRE-GY-9743 -- Assignment 3}
\fancyfoot[C]{\footnotesize\sffamily\color{inkgrey}\thepage}

\setlist[itemize]{leftmargin=1.4em,itemsep=2pt,topsep=3pt,
  label=\textcolor{NYUviolet}{\textbullet}}
\setlist[enumerate]{leftmargin=1.6em,itemsep=2pt,topsep=3pt}
\newcolumntype{L}[1]{>{\raggedright\arraybackslash}p{#1}}
\newcolumntype{Y}{>{\raggedright\arraybackslash}X}
\renewcommand{\arraystretch}{1.12}
\newcommand{\keynote}[1]{\par\medskip\noindent\textbf{\sffamily\color{NYUviolet}Key result.}\enspace #1\par\medskip}
\newcommand{\warningbox}[1]{\par\medskip\noindent\textit{\textbf{Modelling convention.}\enspace #1}\par\medskip}
\newcommand{\focus}[1]{\noindent{\small\itshape\color{inkgrey}
  \textbf{Question focus.} #1}\par\medskip}
\newcommand{\AI}{\operatorname{AI}}
\newcommand{\CTD}{\operatorname{CTD}}
\numberwithin{equation}{section}
\captionsetup{font=small,labelfont={bf,sf,color=NYUviolet},labelsep=period}

\begin{document}
\setlength{\abovedisplayskip}{7pt}
\setlength{\belowdisplayskip}{7pt}
\setlength{\abovedisplayshortskip}{3pt}
\setlength{\belowdisplayshortskip}{4pt}
\setlength{\jot}{2pt}
\hypersetup{pageanchor=false}

\begin{titlepage}
\thispagestyle{empty}
\centering
\vspace*{0.55in}
{\color{NYUviolet}\rule{\textwidth}{2.2pt}}\\[10pt]
{\sffamily\fontsize{16}{19}\selectfont\color{inkgrey}
Fixed Income Analysis Report}\\[14pt]
{\sffamily\bfseries\fontsize{25}{30}\selectfont\color{NYUviolet}
Bond Futures and the\\[3pt]Cheapest-to-Deliver Option}\\[14pt]
{\sffamily\fontsize{15}{18}\selectfont\color{inkgrey}
Forward Carry, Conversion Factors,\\[3pt]
Implied Repo and Scenario Analysis}\\[10pt]
{\color{NYUviolet}\rule{\textwidth}{2.2pt}}\\[28pt]

\begin{minipage}{0.88\textwidth}
\centering\itshape\small\color{inkgrey}
This report solves the three written questions in Assignment 3. It prices the
three-month forward on the specified Treasury, computes the exchange conversion
factor and normalized forward price, recovers the implied repo rate from the
quoted futures price, identifies the cheapest-to-deliver selection rule, and
sets out a practical framework for forecasting CTD switches.
\end{minipage}\\[34pt]

\begin{tabular}{rl}
\sffamily\bfseries Author & Ziyi Yang\\[3pt]
\sffamily\bfseries Course & FRE-GY-9743\\[3pt]
\sffamily\bfseries Source problem & Jianing (Jay) Yao, Assignment 3\\[3pt]
 & September 28, 2026
\end{tabular}

\vfill
{\sffamily\small\color{inkgrey}Assignment 3 \quad\textbullet\quad Bond Futures}\\[2pt]
{\color{NYUviolet}\rule{\textwidth}{1.2pt}}
\end{titlepage}
\hypersetup{pageanchor=true}

%========================================================
% ABSTRACT, KEYWORDS AND CONTENTS
%========================================================
\begin{center}
{\sffamily\large\bfseries\color{NYUviolet}Abstract}
\end{center}
\begin{quote}
\small\itshape
This report analyses the pricing and delivery economics of a Treasury bond
futures contract. It derives the three-month forward value from cash-and-carry,
converts the bond to the exchange's standard 6\% yield basis, and obtains the
normalized forward price. It then recovers the implied repo rate from the
quoted futures price and states the cheapest-to-deliver rule. Finally, it
develops a scenario procedure for forecasting CTD switches as yields, coupon
carry, accrued interest, financing conditions and time to delivery change.
Price units, accrued-interest treatment and financing conventions are stated
explicitly throughout.
\end{quote}
\vspace{4pt}
{\small\noindent\textbf{\sffamily Keywords:}\;
Bond futures \textbullet\; Conversion factor \textbullet\; Implied repo
\textbullet\; Cheapest to deliver \textbullet\; Net basis}
\vspace{8pt}
{\color{rulegrey}\hrule}
\setcounter{tocdepth}{1}
\renewcommand{\contentsname}{\sffamily\large\bfseries\color{NYUviolet}Contents}
{\small\tableofcontents}
{\color{rulegrey}\hrule}
\clearpage

\section*{Notation and Modelling Conventions}
\addcontentsline{toc}{section}{Notation and Modelling Conventions}

All prices are per \$100 face amount. The bond pays a 4\% annual coupon in two
equal installments, so each semiannual coupon is
\[
C=\frac{4}{2}=2.
\]
The statement that the bond has \emph{exactly} nine years remaining is taken to
place the valuation date on a coupon date. Hence the spot accrued interest is
zero. The delivery date is three months later, halfway through the next
six-month coupon period, and no coupon is paid between valuation and delivery.
The repo quote is treated as a simple annual rate with $T=3/12=0.25$.

\begin{table}[H]
\centering
\caption{Inputs and price conventions}
\small
\begin{tabularx}{\textwidth}{@{}L{3.1cm}Y@{}}
\toprule\rowcolor{tablehead}\textbf{Item} & \textbf{Value or convention}\\\midrule
Spot clean price & $S_{0,\mathrm{clean}}=100$\\
Spot accrued interest & $\AI_0=0$, so $S_{0,\mathrm{dirty}}=100$\\
Coupon and frequency & 4\% annually; $C=2$ every six months\\
Remaining maturity & 9 years, or 18 semiannual periods\\
Forward tenor & $T=0.25$ year\\
Market repo rate & $r_{\mathrm{mkt}}=2\%$ simple annual\\
Conversion-factor yield & 6\% nominal, compounded semiannually\\
Delivery accrued interest & $\AI_T=2(3/6)=1$\\
\bottomrule
\end{tabularx}
\end{table}

\keynote{The accrued-interest adjustment is essential. The financed spot
position grows to a dirty value of 100.50, while the comparable quoted forward
price is clean and therefore equals 99.50.}

\clearpage
\section{Q1: Normalized forward price}
\focus{Compute the three-month forward price, the conversion factor, and the
normalized forward price; then explain what the conversion factor accomplishes.}

\subsection{Three-month forward price}
Because no coupon is paid before delivery, cash-and-carry gives the dirty
delivery value
\begin{equation}
B^{\mathrm{dirty}}_F
=S_{0,\mathrm{dirty}}(1+r_{\mathrm{mkt}}T)
=100(1+0.02\times0.25)=100.50.
\end{equation}
At delivery, three months of the next semiannual coupon have accrued:
\begin{equation}
\AI_T=C\frac{3}{6}=2\times\frac12=1.00.
\end{equation}
The clean forward price used with the conversion factor is therefore
\begin{equation}
\boxed{B_F=B^{\mathrm{dirty}}_F-\AI_T=100.50-1.00=99.50.}
\end{equation}

\subsection{Conversion factor}
At the exchange's 6\% standard yield, the semiannual yield is 3\%. With 18
remaining coupon periods, the clean model price is
\begin{align}
P_{6\%}
&=\sum_{i=1}^{18}\frac{2}{(1.03)^i}
  +\frac{100}{(1.03)^{18}}\\
&=86.2464869.
\end{align}
The conversion factor is this clean price divided by 100:
\begin{equation}
\boxed{K=\frac{86.2464869}{100}=0.8624648692.}
\end{equation}

\subsection{Normalized price and interpretation}
The normalized forward price is
\begin{equation}
\boxed{\frac{B_F}{K}
=\frac{99.50}{0.8624648692}
=115.3670179\approx115.3670.}
\end{equation}

The conversion factor maps bonds with different coupons and maturities onto a
common 6\% yield benchmark. A low-coupon bond such as this 4\% issue has
$K<1$, so its clean delivery value is divided by a number below one. This
normalization makes a basket deliverable against one futures contract, but it
is not perfect: the actual yield curve, coupon carry, repo specialness and
accrued interest differ across bonds. Those residual differences create the
cheapest-to-deliver option. Accrued interest is excluded from both clean bond
prices and the conversion factor and is added separately to the delivery
invoice.

\clearpage
\section{Q2: Cheapest-to-deliver and implied repo}
\focus{Use the quoted futures price to recover the bond's implied repo and state
the economic rule for selecting the CTD bond.}

\subsection{Implied repo for the specified bond}
Let the quoted futures price be $F=115.3670$. The clean invoice component is
$FK$ and the dirty amount received upon delivery is $FK+\AI_T$. Since there is
no intervening coupon, the implied simple repo rate solves
\begin{equation}
S_{0,\mathrm{dirty}}(1+r_{\mathrm{imp}}T)=FK+\AI_T.
\end{equation}
Hence
\begin{align}
r_{\mathrm{imp}}
&=\frac{1}{T}
\left(\frac{FK+\AI_T}{S_{0,\mathrm{dirty}}}-1\right)\\
&=\frac{1}{0.25}
\left(\frac{115.3670(0.8624648692)+1}{100}-1\right)\\
&=0.01999938.
\end{align}
Therefore
\begin{equation}
\boxed{r_{\mathrm{imp}}\approx1.99994\%\approx2.00\%.}
\end{equation}
The small difference from exactly 2\% is caused only by rounding the futures
quote to four decimal places. This confirms the consistency between Questions
1 and 2.

\subsection{How to select the CTD bond}
For each deliverable bond $x$, calculate its implied repo rate from the observed
futures price, its conversion factor and all relevant carry cash flows. A trader
can buy the bond, finance it in repo and deliver it into the short futures
position. If its implied repo exceeds its actual financing rate, this trade has
positive carry before costs.

When all bonds can be financed near the same general-collateral rate, the CTD is
the bond with the \textbf{highest implied repo rate}. Equivalently, it has the
lowest normalized forward delivery cost:
\begin{equation}
\boxed{
\CTD=\arg\min_{x\in\mathcal X}\frac{B_{F,x}}{K_x}
=\arg\max_{x\in\mathcal X}r^{\mathrm{imp}}_x.}
\end{equation}
If individual bonds trade at different repo rates, especially when one is
``special,'' the comparison should use financing-adjusted carry. In practice,
the preferred bond maximizes the implied financing advantage
$r^{\mathrm{imp}}_x-r^{\mathrm{actual}}_x$, or equivalently minimizes net basis
after coupon, financing and transaction-cost adjustments.

\keynote{For the bond in Question 1, the quoted futures price 115.3670
produces an implied repo of approximately 2\%, exactly matching the assumed
market repo rate. Its cash-and-carry trade is therefore at break-even before
costs.}

\clearpage
\section{Q3: Scenario analysis and CTD switches}
\focus{Predict which bond may become CTD at a future horizon and explain why the
delivery choice matters.}

\subsection{A reproducible scenario procedure}
At a horizon $h=N/365$, the CTD forecast should be recomputed rather than held
fixed. For every bond $x$ in the delivery basket:
\begin{enumerate}
\item Project its clean price $S_x(h)$ under the scenario yield curve. Roll the
coupon schedule forward and compute the accrued interest at the horizon.
\item Apply the scenario repo rate for that specific bond. Include coupons paid
between $h$ and delivery and their reinvestment or financing treatment.
\item Compute the scenario forward delivery price $B_{F,x}(h)$ and divide by the
exchange-fixed conversion factor $K_x$.
\item Select the minimum normalized forward price. As a cross-check, use the
observed scenario futures price to compute implied repo for every bond and
select the largest financing advantage over actual repo.
\end{enumerate}

Formally, the scenario CTD is
\begin{equation}
x^*(h)=\arg\min_{x\in\mathcal X}
\frac{B_{F,x}\!\left(S_x(h),r_x(h),T_d-h\right)}{K_x}.
\end{equation}

\clearpage
\subsection{What can cause a switch?}
\begin{table}[H]
\centering
\caption{Main drivers of CTD changes}
\small
\begin{tabularx}{\textwidth}{@{}L{3.2cm}Y@{}}
\toprule\rowcolor{tablehead}\textbf{Driver} & \textbf{Effect on the ranking}\\\midrule
Yield level and curve shape & Bonds have different duration and convexity, so a parallel move or curve twist changes their clean prices by different amounts.\\
Coupon carry and accrued interest & Coupon dates, accrued interest and reinvestment alter the cost of carrying each bond to delivery.\\
Repo specialness & A bond that is expensive to borrow or unusually cheap to finance can move sharply in implied-repo and net-basis rankings.\\
Passage of time & Remaining financing periods shorten and coupon entitlements change as the contract approaches delivery.\\
Conversion-factor bias & Conversion factors are fixed from a 6\% benchmark and do not update with the actual yield curve, leaving residual relative-value exposure.\\
\bottomrule
\end{tabularx}
\end{table}

CTD identity matters because the futures price is anchored by the cheapest
delivery route. The CTD determines the relevant cash bond, conversion factor,
duration and DV01 used in a hedge. A CTD switch can therefore change hedge
ratios and produce basis P\&L even when the broad level of rates moves as
expected. For a short futures position it also determines the bond that should
be acquired and delivered; for a long position it describes the quality option
granted to the short.

\keynote{A CTD forecast is a relative-value calculation across the entire
basket. Reprice every candidate under the same market scenario, include
bond-specific financing, and choose the minimum normalized delivery cost rather
than extrapolating today's CTD mechanically.}

\clearpage
\section*{Summary of the Three Answers}
\addcontentsline{toc}{section}{Summary of the Three Answers}
\begin{table}[H]
\centering
\caption{Numerical results and decision rules}
\small
\begin{tabularx}{\textwidth}{@{}L{3.4cm}Y@{}}
\toprule\rowcolor{tablehead}\textbf{Quantity} & \textbf{Answer}\\\midrule
Three-month dirty forward value & $100.5000$\\
Delivery accrued interest & $1.0000$\\
Clean forward price $B_F$ & $99.5000$\\
Conversion factor $K$ & $0.8624648692$\\
Normalized forward price $B_F/K$ & $115.3670179\approx115.3670$\\
Implied repo at $F=115.3670$ & $1.99994\%\approx2.00\%$\\
CTD rule under common GC financing & Highest implied repo, equivalently lowest normalized forward cost\\
CTD rule with bond-specific repo & Highest financing-adjusted implied-repo advantage, equivalently lowest net basis\\
\bottomrule
\end{tabularx}
\end{table}

\end{document}
